% ============================================================
%  glossary.tex: Glossareintraege und Abkuerzungen.
%  Aufruf mit \gls{...}, \glspl{...}, \Gls{...},
%  oder \acrshort{...} / \acrlong{...} / \acrfull{...}.
%  Gedruckt werden nur tatsaechlich zitierte Eintraege.
% ============================================================

\newglossaryentry{altholz}{
  name=Altholz,
  description={Gebrauchtes Holz aus Rueckbau, Abbruch oder Demontage, das
               nach der Altholzverordnung in vier Kategorien eingeteilt wird;
               nur die gering belasteten Kategorien kommen fuer eine tragende
               Wiederverwendung in Betracht}
}

\newglossaryentry{druckzone}{
  name=Druckzone,
  description={Der Querschnittsbereich eines auf Biegung beanspruchten
               Traegers oberhalb der neutralen Faser, in dem die
               Laengsspannungen Druckspannungen sind}
}

\newglossaryentry{zugzone}{
  name=Zugzone,
  description={Der Querschnittsbereich eines auf Biegung beanspruchten
               Traegers unterhalb der neutralen Faser, in dem die
               Laengsspannungen Zugspannungen sind}
}

\newglossaryentry{neutralefaser}{
  name=neutrale Faser,
  description={Die spannungsfreie Schicht im Querschnitt eines gebogenen
               Traegers, die Druck- und Zugzone voneinander trennt}
}

\newglossaryentry{lastgeschichte}{
  name=Lastgeschichte,
  description={Die gesamte Abfolge der Beanspruchungen, die ein Bauteil
               waehrend seiner Nutzungsdauer erfahren hat, einschliesslich
               Dauer, Hoehe und Wechsel der Lasten sowie der begleitenden
               Feuchteverhaeltnisse}
}

\newglossaryentry{kriechen}{
  name=Kriechen,
  description={Die zeitabhaengige, mit der Belastungsdauer zunehmende
               Verformung unter konstanter Last; bei Holz durch Feuchtewechsel
               erheblich verstaerkt (mechano-sorptives Kriechen)}
}

\newglossaryentry{restfestigkeit}{
  name=Restfestigkeit,
  description={Die nach einer Nutzungsphase verbliebene Tragfaehigkeit eines
               Bauteils, im Unterschied zur Festigkeit des ungenutzten
               Ausgangsmaterials}
}

\newglossaryentry{rohdichte}{
  name=Rohdichte,
  description={Masse je Volumeneinheit des Holzes bei definierter Holzfeuchte;
               wichtigster Einzelindikator fuer die mechanischen Eigenschaften
               und daher als Kovariate im Zonenvergleich zu fuehren}
}

\newglossaryentry{sortierung}{
  name=Sortierung,
  description={Die Zuordnung von Bauholz zu einer Festigkeitsklasse, entweder
               visuell nach DIN 4074 oder maschinell nach EN 14081}
}

% --- Abkuerzungen
\newacronym{altholzv}{AltholzV}{Altholzverordnung}
\newacronym{krwg}{KrWG}{Kreislaufwirtschaftsgesetz}
\newacronym{hnee}{HNEE}{Hochschule fuer nachhaltige Entwicklung Eberswalde}
\newacronym{mor}{MOR}{Biegefestigkeit (modulus of rupture)}
\newacronym{moe}{MOE}{Elastizitaetsmodul (modulus of elasticity)}
\newacronym{ndt}{NDT}{zerstoerungsfreie Pruefung (non-destructive testing)}
\newacronym{dol}{DOL}{Lastdauereffekt (duration of load)}
\newacronym{anova}{ANOVA}{Varianzanalyse (analysis of variance)}
